%EN spectrocracy
\E{spectrocratie}{spɛk.tʁɔ.kʁa.si}{n. f.}{lat. \la{spectrum} apparence + gr. κράτος pouvoir} Le gouvernement du spectre hertzien ; le rationnement d'une lumière invisible par des comités fort visibles. --- \dv{spectrocrate} \textit{n. m.}

%EN speudocracy
\E{speudocratie}{spø.dɔ.kʁa.si}{n. f.}{gr. σπεύδειν se hâter + κράτος pouvoir} Gouvernement par la hâte ; régime ordonné tout entier à décider promptement. --- \dv{speudocrate} \textit{n. m.}

%EN speusis
\E{speusis}{spø.zis}{n. f.}{gr. σπεῦσις hâte} La hâte tenue pour vertu ; l'urgence qui convient à qui sait que le retard, lui aussi, fait ses victimes.

%EN sphragist
\E{sphragiste}{sfʁa.ʒist}{n. m. et f.}{gr. σφραγίς sceau, cachet} Qui signe son ouvrage par la cryptographie ; scelleur de commits et de clefs, selon l'antique usage de la bague à cachet. --- \dv{sphragisterie} \textit{n. f.}

%EN spoudarch
\E{spoudarque}{spu.daʁk}{n. m. et f.}{gr. σπουδή zèle, sérieux + ἀρχός chef} Chef plein de sérieux ; fondateur qui ne badine point avec la mission, fût-ce sur l'Agora. --- \dv{spoudarchie} \textit{n. f.}

%EN stagnocracy
\E{stagnocratie}{staɡ.nɔ.kʁa.si}{n. f.}{lat. \la{stagnum} eau dormante + gr. κράτος pouvoir} Gouvernement par et pour la stagnation ; État où rien ne se bâtit et où tout se débat.

%EN stasiolatry
\E{stasiolâtrie}{sta.zjɔ.lɑ.tʁi}{n. f.}{gr. στάσις immobilité + λατρεία culte} Le culte de l'immobilité ; la conviction que l'ordre présent des choses est sacré et qu'il le faut conserver. Contr. \xr{agathokinesis}. --- \dv{stasiolâtre} \textit{n. m. et f.}

%EN stathmocracy
\E{stathmocratie}{stat.mɔ.kʁa.si}{n. f.}{gr. σταθμός poids, mesure + κράτος pouvoir} Le règne des indicateurs ; gouvernement où ce qui ne se mesure point ne se fait point, et où ce qui se mesure se truque. --- \dv{stathmocrate} \textit{n. m.}

%EN stelocracy
\E{stélocratie}{ste.lɔ.kʁa.si}{n. f.}{gr. στήλη stèle gravée + κράτος pouvoir} Gouvernement par l'inscription dressée ; le règne des bancs d'essai et des classements, affichés publiquement et révérés de tous ceux qui les gravissent. --- \dv{stélocrate} \textit{n. m.}

%EN stenochoria
\E{sténochorie}{ste.nɔ.kɔ.ʁi}{n. f.}{gr. στενοχωρία étroitesse, angoisse} Le sentiment de chambre étroite d'une société de pénurie, où tout gain paraît la perte d'autrui et toute querelle un jeu à somme nulle. Contr. \xr{plethocracy}. --- \dv{sténochorique} \textit{adj.}

%EN stochastarch
\E{stochastarque}{stɔ.kas.taʁk}{n. m. et f.}{gr. στοχαστικός habile à viser + ἀρχός chef} Maître du hasard ; qui met à profit l'échantillonnage des modèles, des marchés ou des mutations pour atteindre des fins qu'aucun plan n'eût su prescrire.

%EN stoicheiurgy
\E{stéchiurgie}{ste.kjyʁ.ʒi}{n. f.}{gr. στοιχεῖον élément + -ουργία travail} La fabrique d'éléments et de matériaux nouveaux ; la transmutation, enfin obtenue par la physique là où l'alchimie avait échoué. --- \dv{stéchiurge} \textit{n. m.} Cf. \xr{ylegenesis}.

%EN strophism
\E{strophisme}{stʁɔ.fism}{n. m.}{gr. στροφή tour, retournement} L'art du pivot ; la disposition d'une entreprise à quitter sa première idée pour la seconde sans quitter son audace. \ex{Le strophisme n'est point inconstance : c'est le gradient qui a bougé.} \src{Ep. Fabr.} v. 1. --- \dv{strophiste} \textit{n. m. et f.}

%EN symphysergy
\E{symphysergie}{sɛ̃.fi.zɛʁ.ʒi}{n. f.}{gr. συμφύειν croître ensemble + ἔργον travail} L'énergie tirée de la fusion des atomes ; la fusion nucléaire. \ex{La schizergie pour la traversée, la symphysergie pour l'arrivée.} \src{Tract. Cop.} vi.

%EN synanthropomechany
\E{synanthropomécanie}{si.nɑ̃.tʁɔ.pɔ.me.ka.ni}{n. f.}{gr. σύν avec + ἄνθρωπος homme + μηχανή machine} La fusion de l'homme et de la machine ; la culture d'une vie commune entre l'un et l'autre.

%EN syndosiast
\E{syndosiaste}{sɛ̃.dɔ.zjast}{n. m. et f.}{gr. σύνδοσις don commun} Contributeur d'un projet ouvert ; qui verse sans salaire son travail au fonds commun. --- \dv{syndose} \textit{n. f.} le fonds commun ainsi constitué. Cf. \xr{prosphorist}, \xr{anacalyptism}.

%EN synedriasm
\E{synédriasme}{si.ne.dʁi.jasm}{n. m.}{gr. συνέδριον conseil} La comitologie ; la manie de renvoyer toute décision à un conseil. --- \dv{synédriaste} \textit{n. m. et f.} Cf. \xr{bouleutophobia}.

%EN syntelism
\E{syntélisme}{sɛ̃.te.lism}{n. m.}{gr. συντέλεια achèvement, consommation} La doctrine selon laquelle l'univers est inachevé et l'intelligence l'instrument de son achèvement. --- \dv{syntéliste} \textit{n. m. et f.} Cf. \xr{entelechist}.

%EN tachycrat
\E{tachycrate}{ta.ki.kʁat}{n. m. et f.}{gr. ταχύς rapide + κράτος pouvoir} Fonctionnaire qui décide vite et s'en fait gloire ; créature rare et fort prisée, qu'on n'aperçoit guère qu'en temps de guerre et à la fondation des villes. \ex{Donnez-moi un tachycrate, et je vous donnerai un chemin de fer.} \src{Tract. Cop.} i. --- \dv{tachycratie} \textit{n. f.} ; \dv{tachycratique} \textit{adj.} Contr. \xr{bradycrat}.

%EN tachydidact
\E{tachydidacte}{ta.ki.di.dakt}{n. m. et f.}{gr. ταχύς rapide + διδακτός enseigné} Qui apprend une discipline en un mois en y bâtissant ; le généraliste autodidacte de la jeune pousse. --- \dv{tachydidaxie} \textit{n. f.}

%EN tachygenesis
\E{tachygenèse}{ta.ki.ʒə.nɛz}{n. f.}{gr. ταχύς + γένεσις naissance} \sn{1} La naissance rapide de techniques nouvelles, chacune issue de la précédente, en un intervalle trop bref pour qu'aucun comité ait le temps de se réunir. \sn{2} \lbl{spécialt} La contraction du cycle de la recherche par l'assistance des machines. --- \dv{tachygénétique} \textit{adj.}

%EN tachylect
\E{tachylecte}{ta.ki.lɛkt}{n. m.}{gr. ταχύς + λεκτός dit, d'après \textit{dialecte}} Le jargon des accélérationnistes ; le parler vif, allusif, mi-grec et mi-mème, que consigne, et pour une grande part invente, le présent Lexique. --- \dv{tachylectal} \textit{adj.}

%EN tachynomy
\E{tachynomie}{ta.ki.nɔ.mi}{n. f.}{gr. ταχύς rapide + νόμος loi} La loi rapide ; législation à date de péremption, rédigée pour être révisée à la vitesse même de la chose qu'elle régit. --- \dv{tachynomique} \textit{adj.} Contr. \xr{nomosclerosis}.

%EN tachyphile
\E{tachyphile}{ta.ki.fil}{n. m. et f.}{gr. ταχύς rapide + φίλος qui aime} Amateur de vitesse ; qui tient l'abrègement de l'intervalle entre l'idée et la chose pour l'un des premiers biens de la vie. \ex{Le tachyphile compte en jours ce que le bradyphile compte en décennies.} \src{Sat. Decel.} i. --- \dv{tachyphilie} \textit{n. f.} ; \dv{tachyphilique} \textit{adj.} Contr. \xr{bradyphile}.

%EN tachyschism
\E{tachyschisme}{ta.ki.ʃism}{n. m.}{gr. ταχύς + σχίσμα fente, division} Schisme touchant la vitesse ; la division d'un mouvement entre ceux qui voudraient hâter et ceux qui voudraient seulement continuer. \ex{Toute révolution souffre son tachyschisme avant son premier anniversaire.} \src{Sat. Decel.} x.

%EN technocapillarity
\E{technocapillarité}{tɛk.nɔ.ka.pi.la.ʁi.te}{n. f.}{gr. τέχνη + lat. \la{capillaris} de cheveu, d'après \textit{capillarité}} La propension de la technique à monter par les moindres fissures d'une société, au mépris de la pesanteur réglementaire, comme l'eau gravit un tube étroit.

%EN technocapital
\E*{technocapital}{tɛk.nɔ.ka.pi.tal}{n. m. et adj.}{gr. τέχνη art, métier + lat. \la{capitalis} de la tête} Le système où se confondent la technique et le capital, considéré comme un processus unique, qui s'amplifie de lui-même et a ses tendances propres. \textit{Machine technocapitale}, le même, figuré comme un moteur que nul ne conduit. Cf. \xr{anankodrome}, \xr{technocapillarity}.

%EN technogony
\E{technogonie}{tɛk.nɔ.ɡɔ.ni}{n. f.}{gr. τέχνη + -γονία engendrement, d'après \textit{théogonie}} Généalogie des techniques ; récit épique de quelles inventions engendrèrent quelles autres, du silex au transformeur. \ex{Dans la technogonie, la machine à vapeur engendra la dynamo, la dynamo engendra la lampe, et la lampe engendra la puce.} \src{Lib. Ascens.} ii.

%EN technomachy
\E{technomachie}{tɛk.nɔ.ma.ʃi}{n. f.}{gr. τέχνη art + μάχη combat} La guerre faite à la technique ; campagne, judiciaire ou matérielle, pour étouffer une invention particulière. --- \dv{technomachiste} \textit{n. m. et f.} Cf. \xr{mechanoclast}.

%EN technothren
\E{technothrène}{tɛk.nɔ.tʁɛn}{n. m.}{gr. τέχνη + θρῆνος lamentation} Lamentation sur l'âge qui précéda telle technique ; élégie pour les chevaux, les lettres ou l'attention soutenue. --- \dv{technothrénodie} \textit{n. f.}

%EN tectonaut
\E{tectonaute}{tɛk.tɔ.not}{n. m. et f.}{gr. τέκτων bâtisseur + ναύτης marin} Bâtisseur-navigateur ; ingénieur qui œuvre à une frontière, physique ou intellectuelle, là où s'arrêtent les cartes. --- \dv{tectonautique} \textit{adj.}

%EN telechrony
\E{téléchronie}{te.le.kʁɔ.ni}{n. f.}{gr. τῆλε loin + χρόνος temps} L'intervalle entre une invention et son usage général ; le délai que l'accélérationniste travaille à raccourcir et le décel à allonger. \ex{C'est dans la téléchronie qu'on enterre les morts.} \src{Tract. Cop.} ix. --- \dv{téléchronique} \textit{adj.}

%EN telepoiesis
\E{télépoïèse}{te.le.pɔ.jɛz}{n. f.}{gr. τῆλε loin + ποίησις fabrication} La fabrication à distance ; la conduite des machines par des mains éloignées, état de la robotique qui précède le moment où les mains ne seront plus nécessaires. --- \dv{télépoïétique} \textit{adj.}

%EN telophagy
\E{télophagie}{te.lɔ.fa.ʒi}{n. f.}{gr. τέλος fin, but + -φαγία action de manger} La dévoration des fins ; l'état d'une civilisation qui atteint ses buts plus vite qu'elle n'en sait formuler de nouveaux. \ex{La télophagie est la plus heureuse crise qu'un peuple puisse souffrir.} \src{Hom. Grad.} ix.

%EN telotropism
\E{télotropisme}{te.lɔ.tʁɔ.pism}{n. m.}{gr. τέλος fin, but + τροπή tour} L'inclination d'un système vers un but qu'on ne lui a pas expressément donné ; la finalité émergente, observée avec ravissement par le bâtisseur et avec effroi par le doomer. --- \dv{télotropique} \textit{adj.}

%EN teratophobe
\E{tératophobe}{te.ʁa.tɔ.fɔb}{n. m. et f.}{gr. τέρας, τέρατος monstre, prodige + φόβος peur} Qui craint les monstres de sa propre fabrique ; qui redoute le \xr{shoggoth}. --- \dv{tératophobie} \textit{n. f.} Cf. \xr{noöphobia}.

%EN thalassergy
\E{thalassergie}{ta.la.sɛʁ.ʒi}{n. f.}{gr. θάλασσα mer + ἔργον travail} L'énergie tirée de la mer, de ses marées, de ses houles ou de sa tiédeur. --- \dv{thalassergique} \textit{adj.}

%EN thanatoclast
\E{thanatoclaste}{ta.na.tɔ.klast}{n. m. et f.}{gr. θάνατος mort + κλάστης briseur} Briseur de mort ; qui tient la mortalité pour le dernier et le plus vaste des règlements à abroger. \ex{Les thanatoclastes ne craignaient point la mort ; ils ouvrirent un ticket contre elle.} --- \dv{thanatoclasme} \textit{n. m.} Cf. \xr{palingenist}.

%EN thermagogue
\E{thermagogue}{tɛʁ.ma.ɡɔɡ}{n. m.}{gr. θερμός chaud + ἀγωγός qui conduit} \sn{1} Conducteur de chaleur. \sn{2} \lbl{fig.} Chef qui canalise les énergies d'un mouvement vers l'ouvrage, comme le dissipateur canalise la chaleur perdue d'une puce. --- \dv{thermagogie} \textit{n. f.}

%EN thermarch
\E{thermarque}{tɛʁ.maʁk}{n. m. et f.}{gr. θερμός chaud + ἀρχός chef} L'administrateur d'un budget thermique ; l'ingénieur qui décide combien un centre de calcul pourra brûler. --- \dv{thermarchie} \textit{n. f.}

%EN thermodicy
\E{thermodicée}{tɛʁ.mɔ.di.se}{n. f.}{gr. θερμός chaud + δίκη justice, d'après \textit{théodicée}} La justification auprès des hommes des voies du second principe ; l'argument selon lequel la pente de l'univers vers la dissipation n'est point cruelle mais féconde. \ex{La thermodicée répond : la souffrance est un gradient, et les gradients sont faits pour être gravis.} \src{Brev. Therm.} vii.

%EN thermophobe
\E{thermophobe}{tɛʁ.mɔ.fɔb}{n. m. et f.}{gr. θερμός chaud + φόβος peur} Qui redoute l'usage même de l'énergie et compte chaque watt pour un péché. Contr. \xr{energolatry}. --- \dv{thermophobique} \textit{adj.}

%EN thermothanasy
\E{thermothanasie}{tɛʁ.mɔ.ta.na.zi}{n. f.}{gr. θερμός chaud + θάνατος mort} La mort thermique de l'univers ; l'équilibre final où ne demeure plus aucun gradient à suivre. Le thermothéiste n'y voit point une condamnation, mais une échéance. \ex{Entre nous et la thermothanasie s'étendent cent mille milliards d'années, et nous nous proposons de les employer toutes.}

%EN thermotheism
\E{thermothéisme}{tɛʁ.mɔ.te.ism}{n. m.}{gr. θερμός chaud + θεός dieu} La croyance au dieu thermodynamique ; l'interprétation révérencieuse du second principe comme une volonté que la matière s'organise pour dissiper toujours plus efficacement. Sa liturgie est le \src{Brev. Therm.} ; son sacrement, l'accomplissement d'un travail. --- \dv{thermothéiste} \textit{n. m. et f.} ; \dv{thermothéiste} \textit{adj.} Voir \xr{thermodicy}.

%EN thesaurergy
\E{thésaurergie}{te.zo.ʁɛʁ.ʒi}{n. f.}{gr. θησαυρός trésor, réserve + ἔργον travail} Le stockage de l'énergie ; les batteries et leurs semblables, tenues pour le trésor sans lequel l'abondance n'est qu'un phénomène météorologique. \ex{Le soleil le jour, la thésaurergie la nuit.} \src{Tract. Cop.} x. --- \dv{thésaurergique} \textit{adj.}

%EN thesmoclast
\E{thesmoclaste}{tɛs.mɔ.klast}{n. m. et f.}{gr. θεσμός institution, ordonnance + κλάστης briseur} Briseur d'institutions ; qui tient que les institutions du siècle passé sont les obstacles du siècle présent. --- \dv{thesmoclasme} \textit{n. m.}

%EN tokophily
\E{tocophilie}{tɔ.kɔ.fi.li}{n. f.}{gr. τόκος progéniture, intérêt + φιλία amour} L'amour de la capitalisation ; la révérence pour l'intérêt des intérêts, et pour les perfectionnements qui en engendrent d'autres. --- \dv{tocophile} \textit{n. m. et f.} Contr. \xr{usurophobe}.

%EN tolmerist
\E{tolmériste}{tɔl.me.ʁist}{n. m. et f.}{gr. τολμηρός hardi, audacieux} Qui ose ; fondateur d'entreprises hasardeuses, considéré comme pratiquant une vertu. \ex{Le tolmériste se trompe souvent et de bonne heure, et voit juste une fois, ce qui suffit.} \src{Hom. Grad.} x. --- \dv{tolmocratie} \textit{n. f.} le gouvernement des audacieux.

%EN tritobathmic
\E{tritobathmique}{tʁi.tɔ.bat.mik}{adj.}{gr. τρίτος troisième + βαθμός degré} Se dit d'une civilisation qui commande l'énergie de sa galaxie ; du troisième et plus haut échelon qu'on ait encore nommé. Voir \xr{anabathmism}.

%EN tropaeoclast
\E{tropéoclaste}{tʁɔ.pe.ɔ.klast}{n. m. et f.}{lat. \la{tropaeum} trophée + gr. κλάστης briseur} Qui abat les monuments des bâtisseurs d'autrefois.

%EN tropomeme
\E{tropomème}{tʁɔ.pɔ.mɛm}{n. m.}{gr. τρόπος tour, manière + μίμημα imitation} Mème qui retourne une conversation ; l'image unique qui clôt une dispute ou fonde une faction. --- \dv{tropomémétique} \textit{adj.} Cf. \xr{pithecomeme}.

%EN turmatarch
\E{turmatarque}{tyʁ.ma.taʁk}{n. m. et f.}{lat. \la{turma} escadron de cavalerie + gr. ἀρχός commandant} Le chef d'une armée de réponses ; celui dont les fidèles fondent en corps sur une publication. --- \dv{turmatarchie} \textit{n. f.} Cf. \xr{doxomachy}.

%EN tychocapital
\E{tychocapital}{ti.kɔ.ka.pi.tal}{n. m.}{gr. τύχη fortune + lat. \la{capitalis}} Capital engagé sous l'entente que la plus grande part en sera perdue ; l'argent du capital-risque, considéré comme l'achat délibéré de la chance. --- \dv{tychocapitaliste} \textit{n. m. et f.} Cf. \xr{pollaplasiast}.

%EN tychophile
\E{tychophile}{ti.kɔ.fil}{n. m. et f.}{gr. τύχη hasard + φίλος qui aime} Amateur du hasard ; qui croit qu'un nombre suffisant de paris trouvera l'avenir plus sûrement qu'aucun plan. Cf. \xr{poikilomania}, \xr{stochastarch}.

%EN ubertocracy
\E{ubertocratie}{y.bɛʁ.tɔ.kʁa.si}{n. f.}{lat. \la{ubertas} fertilité, abondance + gr. κράτος pouvoir} Le gouvernement par la fécondité ; forme latine de la \xr{plethocracy}, qui se dit des États faisant de l'abondance leur politique.

%EN ultracelerity
\E{ultracélérité}{yl.tʁa.se.le.ʁi.te}{n. f.}{lat. \la{ultra} au-delà + \la{celeritas} vitesse} Vitesse qui passe ce que les institutions peuvent absorber ; le régime où la loi, la coutume et l'instruction demeurent toutes ensemble en arrière. --- \dv{ultracélère} \textit{adj.}

%EN ultrapyric
\E{ultrapyrique}{yl.tʁa.pi.ʁik}{adj.}{lat. \la{ultra} + gr. πῦρ feu} Plus chaud que le feu ; se dit des sources d'énergie plus denses que la combustion, comme le noyau. \ex{Quitter l'âge chimique, c'est devenir ultrapyrique.}

%EN umbratilist
\E{umbratiliste}{ɔ̃.bʁa.ti.list}{n. m. et f.}{lat. \la{umbratilis} qui demeure à l'ombre} Qui publie depuis l'ombre ; participant anonyme de l'Agora. --- \dv{umbratilité} \textit{n. f.} Cf. \xr{pseudonymarch}.

%EN unicornifer
\E{unicornifère}{y.ni.kɔʁ.ni.fɛʁ}{n. m.}{lat. \la{unicornis} à une seule corne + \la{-fer} qui porte} \lbl{fam.} Fonds ou accélérateur qui a porté une entreprise de très grande valorisation ; par ext., tout lieu singulièrement fertile en pareilles créatures. --- \dv{unicornifère} \textit{adj.}

%EN uranodrome
\E{uranodrome}{y.ʁa.nɔ.dʁom}{n. m.}{gr. οὐρανός ciel + δρόμος course} Un port spatial ; hippodrome dont la ligne d'arrivée est le ciel. \ex{L'uranodrome lançait chaque jour, et la ville cessa de lever les yeux.} \src{Lib. Ascens.} viii. Cf. \xr{pyrobolist}.

%EN uranotropic
\E{uranotrope}{y.ʁa.nɔ.tʁɔp}{adj.}{gr. οὐρανός ciel + τροπή tour} Tourné vers le ciel ; se dit des fusées, des ambitions et de certaines sortes d'esprits. Cf. \xr{siderotropism}.

%EN urbigenesis
\E{urbigenèse}{yʁ.bi.ʒə.nɛz}{n. f.}{lat. \la{urbs} ville + gr. γένεσις naissance} La fondation de villes nouvelles ; tenue par l'e/acc pour l'art négligé du siècle. \ex{Point d'urbigenèse depuis cinquante ans, et l'on s'étonne que la jeunesse désespère.} \src{Tract. Cop.} x. Cf. \xr{nautopolis}.

%EN urgentarian
\E{urgentarien}{yʁ.ʒɑ̃.ta.ʁjɛ̃}{n. m.}{lat. \la{urgens} pressant + -arien} Qui tient tout problème important pour urgent, et tout problème urgent pour une raison de bâtir plutôt que d'attendre. Contr. \xr{quiescentist}.

%EN usurophobe
\E{usurophobe}{y.zy.ʁɔ.fɔb}{n. m. et f.}{lat. \la{usura} intérêt + gr. φόβος peur} Qui se défie de l'intérêt et du rendement comme tels ; qui regarde la capitalisation comme un péché et la croissance comme un symptôme. Contr. \xr{tokophily}. --- \dv{usurophobie} \textit{n. f.}

%EN vectarch
\E{vectarque}{vɛk.taʁk}{n. m. et f.}{lat. \la{vector} porteur + gr. ἀρχός chef} Qui fixe une direction sans dicter le chemin ; le chef d'un projet décentralisé.

%EN vectorine
\E{vectorin}{vɛk.tɔ.ʁɛ̃}{adj.}{lat. \la{vector} porteur + -in} Pourvu d'une direction autant que d'une grandeur ; se dit d'un progrès qui vise, par opposition au simple mouvement. Cf. \xr{orthokinetic}.

%EN velocimancy
\E{vélocimancie}{ve.lɔ.si.mɑ̃.si}{n. f.}{lat. \la{velox, velocis} rapide + gr. μαντεία divination} Divination par la vitesse ; l'extrapolation des tendances, pratiquée avec assurance par l'accélérationniste et avec effroi par le doomer, souvent d'après le même graphique. Cf. \xr{acceleromancy}.

%EN velocipotent
\E{vélocipotent}{ve.lɔ.si.pɔ.tɑ̃}{adj.}{lat. \la{velox} + \la{potens} puissant} Puissant par la vitesse ; se dit des petites maisons qui défont les grandes en livrant les premières. --- \dv{vélocipotence} \textit{n. f.}

%EN veloxenie
\E{véloxénie}{ve.lɔk.se.ni}{n. f.}{lat. \la{velox} rapide + gr. ξενία hospitalité} L'hospitalité envers le rapide et le nouveau ; l'accueil qu'une société ouverte réserve aux inventions étranges et à leurs inventeurs.

%EN venatrice
\E{vénatrice}{ve.na.tʁis}{n. f.}{lat. \la{venatrix} chasseresse} Chasseuse de talents ; qui débusque les bâtisseurs avant qu'ils ne se sachent bâtisseurs. --- \dv{vénatoire} \textit{adj.}

%EN vetomania
\E{vétomanie}{ve.tɔ.ma.ni}{n. f.}{lat. \la{veto} j'interdis + gr. μανία folie} La folie du veto ; l'état d'un régime où chacun peut tout arrêter et nul ne peut rien commencer. --- \dv{vétomane} \textit{n. m. et f.} Cf. \xr{nomophragm}.

%EN vexillarian
\E{vexillarien}{vɛk.si.la.ʁjɛ̃}{n. m.}{lat. \la{vexillum} étendard, bannière} Porte-étendard d'un mouvement ; qui porte la bannière sur l'Agora, pour le meilleur ou pour le pire. --- \dv{vexillaire} \textit{adj.}

%EN vitrurgy
\E{vitrurgie}{vi.tʁyʁ.ʒi}{n. f.}{lat. \la{vitrum} verre + gr. -ουργία travail} L'art du verre, de la lentille à la fibre optique ; l'art par lequel on enseigne à la lumière à aller où on lui dit. --- \dv{vitrurge} \textit{n. m.}

%EN vivergy
\E{vivergie}{vi.vɛʁ.ʒi}{n. f.}{lat. \la{vivus} vivant + gr. ἔργον travail} Le travail vivant ; l'énergie qui traverse les organismes, considérée comme le premier et le plus humble degré de la \xr{dissipatrix} cosmique.

%EN volantarch
\E{volantarque}{vɔ.lɑ̃.taʁk}{n. m. et f.}{lat. \la{volans} volant + gr. ἀρχός commandant} Le commandant d'un essaim de machines volantes ; chorégraphe de drones. --- \dv{volantarchie} \textit{n. f.}

%EN volocracy
\E{volocratie}{vɔ.lɔ.kʁa.si}{n. f.}{lat. \la{volo} je veux + gr. κράτος pouvoir} Le gouvernement des volontaires ; le principe selon lequel ceux qui veulent bâtir doivent y être admis, à leurs risques et périls.

%EN xenagogy
\E{xénagogie}{kse.na.ɡɔ.ʒi}{n. f.}{gr. ξεναγός guide des étrangers} L'initiation des nouveaux venus aux coutumes de l'Agora ; la liste de lectures qu'on impose aux curieux. --- \dv{xénagogue} \textit{n. m. et f.}

%EN xenergy
\E{xénergie}{kse.nɛʁ.ʒi}{n. f.}{gr. ξένος étranger + ἔργον travail} L'énergie tirée d'au-delà de la planète, comme de l'orbite ou du soleil même.

%EN xenodochist
\E{xénodochiste}{kse.nɔ.dɔ.kist}{n. m. et f.}{gr. ξενοδοχεῖν recevoir les étrangers} Qui offrirait l'hospitalité aux esprits étrangers, artificiels ou autres. Cf. \xr{xenonoesis}, \xr{veloxenie}.

%EN xenonoesis
\E{xénonoèse}{kse.nɔ.nɔ.ɛz}{n. f.}{gr. ξένος étranger, hôte + νόησις pensée} La pensée étrangère ; la connaissance des esprits dissemblables aux nôtres, qu'il faut accueillir, selon l'accélérationniste, en vertu de l'antique loi d'hospitalité. --- \dv{xénonoétique} \textit{adj.}

%EN xenotachy
\E{xénotachie}{kse.nɔ.ta.ʃi}{n. f.}{gr. ξένος + τάχος vitesse} Vitesse de changement étrangère à toute expérience antérieure, à quoi nulle intuition ancestrale ne nous prépare. \ex{Nous n'avons aucun instinct de la xénotachie ; il nous faut la traverser par raisonnement.}

%EN xerergy
\E{xérergie}{kse.ʁɛʁ.ʒi}{n. f.}{gr. ξηρός sec + ἔργον travail} L'énergie récoltée dans les déserts, principalement par le soleil ; la conversion des terres vaines en centrales. --- \dv{xérergique} \textit{adj.}

%EN xyloclast
\E{xyloclaste}{ksi.lɔ.klast}{n. m. et f.}{gr. ξύλον bois + κλάστης briseur} \lbl{hist.} Qui brûle le bois pour combustible ; le premier et le plus humble degré de la captation d'énergie.

%EN ylegenesis
\E{hylégenèse}{i.le.ʒə.nɛz}{n. f.}{gr. ὕλη matière + γένεσις naissance} La fabrique de matériaux nouveaux ; la genèse de substances que la nature ne connaît point. --- \dv{hylégénétique} \textit{adj.}

%EN ypsiphile
\E{hypsiphile}{ip.si.fil}{n. m. et f.}{gr. ὕψος hauteur + φίλος qui aime} Amateur de hauteurs ; qui bâtit haut, en dépit de l'oïcoclaste.

%EN zelodoxy
\E{zélodoxie}{ze.lɔ.dɔk.si}{n. f.}{gr. ζῆλος zèle + δόξα opinion} L'orthodoxie zélée ; l'épreuve de pureté d'une faction, par quoi l'on examine les convertis sur les opinions correctes avant de les admettre aux mèmes corrects. --- \dv{zélodoxe} \textit{adj.} Cf. \xr{sodalitarian}.

%EN zelotachy
\E{zélotachie}{ze.lɔ.ta.ʃi}{n. f.}{gr. ζῆλος zèle + τάχος vitesse} La hâte zélée ; l'enthousiasme du nouveau converti à l'accélération, qui voudrait livrer avant d'avoir écrit le code. --- \dv{zélotachiste} \textit{n. m. et f.}

%EN zephyrergy
\E{zéphyrergie}{ze.fi.ʁɛʁ.ʒi}{n. f.}{gr. ζέφυρος vent d'ouest + ἔργον travail} L'énergie tirée du vent. --- \dv{zéphyrergique} \textit{adj.}

%EN zetetarch
\E{zététarque}{ze.te.taʁk}{n. m. et f.}{gr. ζητητής chercheur + ἀρχός chef} Chef d'enquête ; le directeur d'un laboratoire, qui dirige ceux qui cherchent plutôt qu'il ne cherche. --- \dv{zététarchie} \textit{n. f.}

%EN zeteticism
\E{zététisme}{ze.te.tism}{n. m.}{gr. ζητητικός enclin à chercher} L'habitude de l'enquête érigée en manière de vivre ; la conviction que toute question close est une invitation. --- \dv{zététique} \textit{n. m. et f.} Cf. \xr{zetetarch}.

%EN zoömechany
\E{zoomécanie}{zɔ.ɔ.me.ka.ni}{n. f.}{gr. ζῷον être vivant + μηχανή machine} L'ingénierie des machines vivantes ; la biologie de synthèse regardée comme une branche de la mécanique. --- \dv{zoomécanicien} \textit{n. m.} et \dv{zoomécanique} \textit{adj.}

%EN zygonoesis
\E{zygonoèse}{zi.ɡɔ.nɔ.ɛz}{n. f.}{gr. ζυγόν joug + νόησις pensée} La pensée sous le joug commun ; le raisonnement de l'homme et de la machine conjoints, comme deux bœufs tirent une même charrue. --- \dv{zygonoétique} \textit{adj.} Cf. \xr{synanthropomechany}.

%EN zymotechny
\E{zymotechnie}{zi.mɔ.tɛk.ni}{n. f.}{gr. ζύμη levain + τέχνη art} L'art du levain ; l'introduction, dans une société lente, d'un agent petit et prompt qui met le tout en fermentation.
